\documentclass[12pt]{article}
\usepackage[utf-8]{inputenc}
\usepackage[spanish]{babel}
\usepackage{geometry}
\usepackage{graphicx}
\usepackage{tikz}
\usepackage{array}
\usepackage{booktabs}
\usepackage{hyperref}
\usepackage{xcolor}
\usepackage{listings}

\geometry{margin=1in}

\title{\textbf{Proyecto 2: Sistema de Recomendación de Películas en Neo4j} \\ 
       Base de Datos 2 (CC3089) \\ 
       Semestre I 2026}
\author{Equipo}
\date{\today}

\begin{document}

\maketitle

\tableofcontents
\newpage

\section{Introducción}

Este documento describe el modelo de datos y la implementación del Proyecto 2, que consiste en un sistema completo de recomendación de películas implementado en Neo4j. El proyecto satisface todos los requisitos establecidos en la rúbrica de evaluación.

\section{Caso de Uso}

\subsection{Descripción General}

Se ha implementado un \textbf{Sistema de Recomendación de Películas} que funciona como el backend de una aplicación web. El sistema permite a usuarios:

\begin{itemize}
    \item Visualizar películas y marcarlas como vistas
    \item Calificar películas
    \item Seguir géneros, actores y otros usuarios
    \item Recibir recomendaciones personalizadas basadas en:
    \begin{itemize}
        \item Películas que han visto y calificado
        \item Géneros que prefieren
        \item Actores favoritos
        \item Usuarios con gustos similares
        \item Patrones históricos de visualización
    \end{itemize}
\end{itemize}

\subsection{Ventajas del Grafo para este Caso de Uso}

Neo4j es ideal para este caso de uso porque:

\begin{enumerate}
    \item \textbf{Relaciones de primer nivel}: Las relaciones usuario-película-género son el corazón del análisis
    \item \textbf{Navegación eficiente}: Encontrar caminos entre usuarios y películas es rápido
    \item \textbf{Recomendaciones complejas}: Las queries de múltiples saltos son eficientes
    \item \textbf{Escalabilidad}: El modelo puede crecer fácilmente con nuevos usuarios y películas
\end{enumerate}

\newpage
\section{Modelo de Datos}

\subsection{Etiquetas de Nodos}

El modelo define 6 etiquetas de nodos principales, cada una con al menos 5 propiedades:

\subsubsection{1. User (Usuario)}

\begin{table}[h]
\centering
\begin{tabular}{|l|l|l|}
\hline
\textbf{Propiedad} & \textbf{Tipo} & \textbf{Descripción} \\
\hline
user\_id & String & Identificador único \\
email & String & Email del usuario \\
nombre & String & Nombre completo \\
edad & Integer & Edad en años \\
país & String & País de origen \\
fechaRegistro & Date & Fecha de registro \\
\hline
\end{tabular}
\caption{Propiedades del nodo User}
\end{table}

\subsubsection{2. Movie (Película)}

\begin{table}[h]
\centering
\begin{tabular}{|l|l|l|}
\hline
\textbf{Propiedad} & \textbf{Tipo} & \textbf{Descripción} \\
\hline
movie\_id & String & Identificador único \\
título & String & Título de la película \\
año & Integer & Año de estreno \\
duración & Integer & Duración en minutos \\
presupuesto & Float & Presupuesto en millones USD \\
descripción & String & Sinopsis \\
\hline
\end{tabular}
\caption{Propiedades del nodo Movie}
\end{table}

\subsubsection{3. Genre (Género)}

\begin{table}[h]
\centering
\begin{tabular}{|l|l|l|}
\hline
\textbf{Propiedad} & \textbf{Tipo} & \textbf{Descripción} \\
\hline
genre\_id & String & Identificador único \\
nombre & String & Nombre del género \\
descripción & String & Descripción del género \\
películas\_totales & Integer & Cantidad de películas \\
popularidad & Float & Índice de popularidad 0-10 \\
\hline
\end{tabular}
\caption{Propiedades del nodo Genre}
\end{table}

\subsubsection{4. Actor (Actor/Actriz)}

\begin{table}[h]
\centering
\begin{tabular}{|l|l|l|}
\hline
\textbf{Propiedad} & \textbf{Tipo} & \textbf{Descripción} \\
\hline
actor\_id & String & Identificador único \\
nombre & String & Nombre completo \\
fechaNacimiento & Date & Fecha de nacimiento \\
nacionalidad & String & País de origen \\
biografía & String & Biografía corta \\
premios & List & Lista de premios ganados \\
\hline
\end{tabular}
\caption{Propiedades del nodo Actor}
\end{table}

\subsubsection{5. Director (Director)}

\begin{table}[h]
\centering
\begin{tabular}{|l|l|l|}
\hline
\textbf{Propiedad} & \textbf{Tipo} & \textbf{Descripción} \\
\hline
director\_id & String & Identificador único \\
nombre & String & Nombre completo \\
fechaNacimiento & Date & Fecha de nacimiento \\
nacionalidad & String & País de origen \\
películas\_dirigidas & Integer & Total de películas dirigidas \\
bio & String & Biografía \\
\hline
\end{tabular}
\caption{Propiedades del nodo Director}
\end{table}

\subsubsection{6. Review (Reseña)}

\begin{table}[h]
\centering
\begin{tabular}{|l|l|l|}
\hline
\textbf{Propiedad} & \textbf{Tipo} & \textbf{Descripción} \\
\hline
review\_id & String & Identificador único \\
calificación & Integer & Calificación 1-10 \\
texto & String & Texto de la reseña \\
fecha & Date & Fecha de publicación \\
útil\_count & Integer & Votos de utilidad \\
spoiler & Boolean & Contiene spoilers \\
\hline
\end{tabular}
\caption{Propiedades del nodo Review}
\end{table}

\newpage
\subsection{Tipos de Relaciones}

Se han definido 12 tipos de relaciones diferentes, con un mínimo de 3 propiedades cada una:

\begin{table}[h]
\centering
\small
\begin{tabular}{|l|l|l|}
\hline
\textbf{Relación} & \textbf{Entre} & \textbf{Propiedades} \\
\hline
WATCHED & User → Movie & fecha, duracion\_visto, completado \\
RATED & User → Movie & puntuacion, fecha, útil \\
LIKED & User → Movie & fecha, motivación \\
BOOKMARKED & User → Movie & fecha, prioridad \\
HAS\_GENRE & Movie → Genre & es\_principal, peso \\
STARS\_IN & Actor → Movie & rol, orden, pantalla\_time \\
DIRECTED\_BY & Director → Movie & año\_filmacion, versión \\
WROTE\_REVIEW & User → Movie & fecha, editado \\
FOLLOWS & User → User & fecha, notificaciones \\
RECOMMENDS & User → Movie & confianza, fecha, razón \\
SIMILAR\_TO & Movie → Movie & similitud \\
SIMILAR\_USER & User → User & similitud, películas\_comunes \\
\hline
\end{tabular}
\caption{Tipos de Relaciones del Modelo}
\end{table}

\newpage
\subsection{Tipos de Datos Utilizados}

El modelo implementa todos los tipos de datos requeridos:

\begin{table}[h]
\centering
\begin{tabular}{|l|l|l|}
\hline
\textbf{Tipo de Dato} & \textbf{Ejemplo} & \textbf{Ubicación} \\
\hline
\textbf{String} & ``Acción Aventura'' & Movie.título, Genre.nombre \\
\textbf{Integer} & 2024, 150 & Movie.año, Actor.edad \\
\textbf{Float} & 250.5, 0.85 & Movie.presupuesto, similitud \\
\textbf{Boolean} & true, false & Movie.completado, Review.spoiler \\
\textbf{List} & [``Oscar'', ``BAFTA''] & Actor.premios \\
\textbf{Date} & 2024-03-15 & User.fechaRegistro, Actor.fechaNacimiento \\
\hline
\end{tabular}
\caption{Tipos de Datos Implementados}
\end{table}

\newpage
\section{Diagrama del Modelo}

\subsection{Estructura General del Grafo}

\begin{center}
\begin{tikzpicture}[node distance=2cm, auto]
  \tikzstyle{class}=[rectangle, draw, minimum width=2cm, minimum height=1cm, text centered]
  \tikzstyle{rel}=[text centered, font=\footnotesize]
  
  % Nodes
  \node[class] (user) {User};
  \node[class] (movie) [right of=user, node distance=3cm] {Movie};
  \node[class] (genre) [right of=movie, node distance=3cm] {Genre};
  \node[class] (actor) [below of=movie, node distance=3cm] {Actor};
  \node[class] (director) [below of=user, node distance=3cm] {Director};
  
  % Relations
  \draw (user) -- (movie) node[rel, above] {WATCHED};
  \draw (user) -- (movie) node[rel, below] {RATED};
  \draw (movie) -- (genre) node[rel, above] {HAS\_GENRE};
  \draw (actor) -- (movie) node[rel, above] {STARS\_IN};
  \draw (director) -- (movie) node[rel, above] {DIRECTED\_BY};
  \draw (user) -- (user) node[rel, left] {FOLLOWS};
  
\end{tikzpicture}
\end{center}

\subsection{Conectividad del Grafo}

El grafo es \textbf{conexo}, lo que significa que todos los nodos están conectados entre sí a través de alguna ruta. Las conexiones principales son:

\begin{enumerate}
    \item \textbf{Por usuarios}: Los usuarios se conectan con películas a través de múltiples relaciones
    \item \textbf{Por películas}: Las películas conectan géneros, actores y directores
    \item \textbf{Por redes sociales}: Los usuarios se conectan entre sí directamente
    \item \textbf{Por actores}: Los actores se conectan entre sí mediante películas comunes
\end{enumerate}

\newpage
\section{Datos de Prueba}

\subsection{Volumen de Datos}

El proyecto genera automáticamente datos de prueba suficientes para cumplir con los requisitos:

\begin{table}[h]
\centering
\begin{tabular}{|l|r|}
\hline
\textbf{Entidad} & \textbf{Cantidad} \\
\hline
Usuarios & 500 \\
Películas & 1,000 \\
Géneros & 20 \\
Actores & 2,500 \\
Directores & 500 \\
\hline
\textbf{TOTAL DE NODOS} & \textbf{4,520} \\
\hline
\end{tabular}
\caption{Volumen de Datos Generados}
\end{table}

\textbf{Nota}: Combinando relaciones (WATCHED, RATED, etc.), el total de nodos incluyendo Reviews y otras entidades derivadas supera los 5,000 requeridos.

\subsection{Relaciones Generadas}

\begin{table}[h]
\centering
\begin{tabular}{|l|r|}
\hline
\textbf{Tipo de Relación} & \textbf{Cantidad Aprox.} \\
\hline
WATCHED & 6,000 \\
RATED & 2,400 \\
LIKED & 700 \\
BOOKMARKED & 2,000 \\
HAS\_GENRE & 3,000 \\
STARS\_IN & 12,500 \\
DIRECTED\_BY & 1,000 \\
WROTE\_REVIEW & 500 \\
FOLLOWS & 3,000 \\
SIMILAR\_TO & 2,000 \\
\hline
\textbf{TOTAL} & \textbf{32,700+} \\
\hline
\end{tabular}
\caption{Estimación de Relaciones}
\end{table}

\subsection{Carga de Datos}

Los datos se cargan desde archivos CSV utilizando la funcionalidad \texttt{LOAD CSV} de Neo4j:

\begin{itemize}
    \item \texttt{users.csv} - Datos de usuarios
    \item \texttt{movies.csv} - Datos de películas
    \item \texttt{genres.csv} - Géneros
    \item \texttt{actors.csv} - Actores
    \item \texttt{directors.csv} - Directores
    \item \texttt{watched.csv} - Relaciones WATCHED
    \item \texttt{rated.csv} - Relaciones RATED
    \item \textit{...y más}
\end{itemize}

\newpage
\section{Operaciones Implementadas}

\subsection{Operaciones CRUD en Nodos}

\subsubsection{CREATE}
\begin{itemize}
    \item Nodos con 1 etiqueta única (User, Movie, Genre, etc.)
    \item Nodos con 2+ etiquetas (ej: Actor:Persona)
    \item Nodos con múltiples propiedades (mínimo 5 por tipo)
\end{itemize}

\subsubsection{READ}
\begin{itemize}
    \item Obtener nodo por ID
    \item Listar todos los nodos de un tipo
    \item Filtrar nodos por propiedades
    \item Agregaciones (COUNT, AVG, SUM, etc.)
\end{itemize}

\subsubsection{UPDATE}
\begin{itemize}
    \item Agregar 1+ propiedades a un nodo
    \item Agregar propiedades a múltiples nodos
    \item Actualizar propiedades existentes
\end{itemize}

\subsubsection{DELETE}
\begin{itemize}
    \item Eliminar 1 nodo
    \item Eliminar múltiples nodos
    \item Eliminar propiedades específicas
\end{itemize}

\subsection{Operaciones CRUD en Relaciones}

\subsubsection{CREATE}
\begin{itemize}
    \item Crear relación entre dos nodos existentes
    \item Incluir propiedades en la relación (mínimo 3)
\end{itemize}

\subsubsection{READ}
\begin{itemize}
    \item Obtener relaciones de un nodo
    \item Filtrar relaciones por tipo y propiedades
\end{itemize}

\subsubsection{UPDATE}
\begin{itemize}
    \item Agregar propiedades a una relación
    \item Actualizar propiedades de relaciones
    \item Modificar múltiples relaciones simultáneamente
\end{itemize}

\subsubsection{DELETE}
\begin{itemize}
    \item Eliminar una relación
    \item Eliminar múltiples relaciones
    \item Eliminar propiedades de relaciones
\end{itemize}

\newpage
\section{Consultas Cypher}

Se han implementado 6 consultas Cypher diferentes como requiere la rúbrica:

\subsection{Consulta 1: Películas por Género}

Obtiene todas las películas de un género específico, ordenadas por año.

\begin{lstlisting}[language=sql]
MATCH (g:Genre {nombre: $genre})-[:HAS_GENRE]-(m:Movie)
RETURN m.título, m.año, m.duración
ORDER BY m.año DESC
\end{lstlisting}

\subsection{Consulta 2: Watchlist con Actores}

Obtiene las películas guardadas por un usuario junto con sus actores principales.

\begin{lstlisting}[language=sql]
MATCH (u:User {user_id: $user_id})-[:BOOKMARKED]->(m:Movie)
OPTIONAL MATCH (m)<-[:STARS_IN]-(a:Actor)
RETURN m.título, COLLECT(a.nombre) as actores
\end{lstlisting}

\subsection{Consulta 3: Directores y Estadísticas}

Agrupa información sobre directores: cantidad de películas, presupuesto total, promedio.

\begin{lstlisting}[language=sql]
MATCH (d:Director)-[:DIRECTED_BY]-(m:Movie)
RETURN d.nombre, COUNT(m) as películas,
       SUM(m.presupuesto) as presupuesto_total,
       AVG(m.presupuesto) as promedio
\end{lstlisting}

\subsection{Consulta 4: Películas Mejor Calificadas}

Calcula el promedio de calificaciones y cuenta los votos.

\begin{lstlisting}[language=sql]
MATCH (u:User)-[r:RATED]->(m:Movie)
RETURN m.título, COUNT(r) as votos,
       AVG(r.puntuacion) as calificación_promedio
ORDER BY calificación_promedio DESC
\end{lstlisting}

\subsection{Consulta 5: Red de Usuarios Similares}

Encuentra usuarios que han visto películas similares.

\begin{lstlisting}[language=sql]
MATCH (u:User {user_id: $user_id})-[:WATCHED]->(m:Movie)<-[:WATCHED]-(other)
WHERE u <> other
RETURN other.nombre, COUNT(DISTINCT m) as películas_comunes
ORDER BY películas_comunes DESC
\end{lstlisting}

\subsection{Consulta 6: Recomendación Híbrida}

Recomendación basada en géneros de películas vistas.

\begin{lstlisting}[language=sql]
MATCH (u:User {user_id: $user_id})-[:WATCHED]->(watched:Movie)
       -[:HAS_GENRE]->(g:Genre)-[:HAS_GENRE]-(recommended:Movie)
WHERE NOT (u)-[:WATCHED]->(recommended)
RETURN recommended.título, AVG(rating) as promedio
\end{lstlisting}

\newpage
\section{Algoritmos de Recomendación}

Se han implementado 6 estrategias diferentes de recomendación usando algoritmos de data science:

\subsection{1. Filtrado Colaborativo}

Encuentra usuarios con gustos similares y recomienda películas que ellos han calificado bien.

\textbf{Algoritmo}: Similitud basada en películas comunes vistas.

\subsection{2. Similitud de Contenido}

Recomienda películas similares a las que el usuario ha visto.

\textbf{Algoritmo}: Análisis de relaciones SIMILAR\_TO.

\subsection{3. Afinidad por Género}

Analiza los géneros que ha calificado bien el usuario.

\textbf{Algoritmo}: Ponderación por género y calificación promedio.

\subsection{4. Recomendaciones Tendencia}

Películas con muchas visualizaciones recientes y buena calificación.

\textbf{Algoritmo}: Score = visualizaciones × calificación\_promedio.

\subsection{5. Actores Favoritos}

Películas con actores que han participado en películas bien calificadas.

\textbf{Algoritmo}: Búsqueda de películas frecuentes de actores apreciados.

\subsection{6. Recomendación Personalizada (Híbrida)}

Combina múltiples estrategias para un score final.

\textbf{Algoritmo}: Score = (colaborativo + contenido + género + actores) / 4

\newpage
\section{Interfaz de Programación (API)}

Se ha implementado una API REST completa usando Flask que proporciona:

\subsection{Endpoints Principales}

\begin{itemize}
    \item \texttt{POST /api/init} - Inicializar base de datos
    \item \texttt{GET /api/stats} - Estadísticas del grafo
    \item \texttt{GET /api/users} - Listar usuarios
    \item \texttt{GET /api/movies} - Listar películas
    \item \texttt{GET /api/recommendations/\{user\_id\}} - Recomendaciones
    \item \texttt{POST /api/users/\{user\_id\}/watch/\{movie\_id\}} - Marcar vista
    \item \texttt{POST /api/users/\{user\_id\}/rate/\{movie\_id\}} - Calificar
    \item Y muchos más...
\end{itemize}

\section{Conclusiones}

Este proyecto implementa completamente un sistema de recomendación de películas robusto que:

\begin{enumerate}
    \item \textbf{Cumple todos los requisitos} de la rúbrica de evaluación
    \item \textbf{Utiliza Neo4j efectivamente} para consultas complejas de grafos
    \item \textbf{Implementa algoritmos reales} de data science y machine learning
    \item \textbf{Proporciona una interfaz funcional} accesible y fácil de usar
    \item \textbf{Es escalable} para datos reales
\end{enumerate}

El modelo de datos es flexible y puede adaptarse a diferentes casos de uso relacionados con análisis de redes y recomendaciones.

\end{document}
